% Generated draft: requires mathematical and source review.
\documentclass[11pt]{article}
\usepackage[margin=25mm]{geometry}
\usepackage{fontspec}
\setmainfont{Latin Modern Roman}
\usepackage{amsmath,amssymb,mathtools,unicode-math}
\setmathfont{Latin Modern Math}
\usepackage{enumitem,needspace,xurl,hyperref}
\setlength{\parindent}{0pt}
\setlength{\parskip}{5pt}
\setlength{\emergencystretch}{3em}
\newcommand{\N}{\mathbb{N}}
\newcommand{\Z}{\mathbb{Z}}
\newcommand{\Q}{\mathbb{Q}}
\newcommand{\R}{\mathbb{R}}
\newcommand{\C}{\mathbb{C}}
\newcommand{\F}{\mathbb{F}}
\newcommand{\A}{\mathbb{A}}
\newcommand{\PP}{\mathbb{P}}
\newcommand{\E}{\mathbb{E}}
\newcommand{\eps}{\varepsilon}
\newcommand{\dd}{\,\mathrm{d}}
\begin{document}
\section*{UnsolvedMath ID: 700775}
\textbf{Source problem code:} catalogue entry
\subsection{Definitions and mathematical statement}
A group \(G\) is \emph{abelian-by-polycyclic} if it has an abelian normal subgroup \(A\trianglelefteq G\) such that \(G/A\) is polycyclic. A group \(P\) is polycyclic if it admits a finite subnormal series
\[
1=P_0\trianglelefteq P_1\trianglelefteq\cdots\trianglelefteq P_n=P
\]
whose factors \(P_i/P_{i-1}\) are cyclic, either finite or infinite. No finite generation assumption is imposed on \(A\).

For a finite generating set \(S\) of \(G\), let \((S\cup S^{-1})^*\) denote the finite words in the generators and their formal inverses, and write \(\overline{w}\in G\) for the element represented by a word \(w\). The conjugacy problem for \(G\) is \emph{soluble} if there is an algorithm which, on every input pair \(u,v\in(S\cup S^{-1})^*\), halts and decides whether
\[
\exists x\in G\qquad x^{-1}\overline{u}x=\overline{v}.
\]
Solubility in this sense is independent of the choice of finite generating set.

\textbf{Question (V. N. Remeslennikov, 1990).} Does every finitely generated abelian-by-polycyclic group have a soluble conjugacy problem?

Here the algorithm may depend on the fixed group and its finite generating set. The supplied statement does not specify a uniform algorithm taking a description of the group as additional input, or an effective format for such descriptions; that stronger interpretation is left unresolved.
\subsection{Short English statement}
Must every finitely generated group with an abelian normal subgroup and polycyclic quotient admit an algorithm deciding whether any two words represent conjugate elements? The question uses the usual interpretation for each fixed group; uniformity over descriptions of groups is unspecified.
\subsection{Sources}
\begin{itemize}
\item Local UnsolvedMath registry, supplied complete source record: permanent ID \(700775\), source problem code \texttt{catalogue entry}, title \emph{Kourovka Notebook Problem 11.83}. The record attributes the problem to V. N. Remeslennikov and gives the proposed year as 1990.
\item Supplied source attribution: \emph{The Kourovka Notebook (Unsolved Problems in Group Theory)}, 11th issue (1990), Problem 11.83; also identified in the record with the 21st issue edition (2026), arXiv identifier \texttt{1401.0300v48}. The collection description names E. I. Khukhro and V. D. Mazurov as editors.
\item The local record labels the problem \emph{open}. This is a supplied registry classification, not an independently verified assessment of the literature. The cited source text and edition have not been independently consulted for this draft.
\end{itemize}
\end{document}
